% Carte des dependances du memoire: chapitres de l'etat des lieux (colonne de
% gauche) vers les chapitres de contributions et de perspectives (colonne de
% droite). Chaque arete porte les sections effectivement mobilisees; la couleur
% designe la destination. Utilise dans FrontBackmatter/CommentLire.tex.
%
% Revision du 2026-08-27: insertion du \cref{ch:learning} entre le chapitre de
% detection et celui sur SPDnet. Il mobilise trois des quatre chapitres de
% l'etat des lieux --- pas le chapitre 4, ce qui laisse au \cref{ch:spdnet} son
% statut de seul chapitre a les mobiliser tous les quatre. Les deux notions
% nouvelles dont il a besoin (modeles graphiques, matrices aleatoires) sont
% introduites dans le chapitre lui-meme et n'apparaissent donc pas ici.
%
% Revision du 2026-08-24: le \cref{ch:spdnet} etant ecrit, ses dependances
% reelles sont desormais connues et elles sont bien plus larges que la seule
% arete depuis le chapitre 3 qui figurait ici. Il mobilise en fait les quatre
% chapitres de l'etat des lieux --- c'est meme le seul chapitre du memoire dans
% ce cas. Trois aretes ont donc ete ajoutees et une entree (§2.1.1, Wishart) a
% ete ouverte dans le bloc du chapitre 2.
\definecolor{lctx}{RGB}{90,90,90}
\definecolor{lcinq}{RGB}{40,70,140}
\definecolor{lsix}{RGB}{110,80,140}
\definecolor{lsept}{RGB}{50,110,120}
\definecolor{lhuit}{RGB}{125,125,125}
\begin{tikzpicture}[
    >={Stealth[length=4pt,width=3pt]},
    line join=round,
    font=\footnotesize,
    entry/.style={anchor=west, inner sep=1pt, font=\scriptsize, text depth=0.5ex},
    ctitle/.style={anchor=west, inner sep=1pt, font=\scriptsize\bfseries, text=lctx},
    frame/.style={draw=lctx!45, rounded corners=2pt, inner xsep=5pt, inner ysep=6pt},
    dest/.style={draw, rounded corners=2pt, align=center, font=\scriptsize,
                 text width=2.5cm, inner sep=5pt},
    elab/.style={font=\tiny, fill=white, inner sep=1.5pt, align=center},
    edge/.style={thin, ->},
  ]

  %--------------------------------------------------------------- etat des lieux
  % Les numeros de chapitre repartent a 1 dans chaque partie
  % (\counterwithin*{chapter}{part}): on les prefixe donc par la partie.
  \node[font=\scriptsize\bfseries, text=lctx, anchor=west] at (0,1.35)
    {Partie \textsc{ii} --- État des lieux};
  \node[ctitle] (t1) at (0,0.62) {ch.\,1 --- Covariance};
  \node[entry]  (a1) at (0, 0.00) {\S1.3 \ estimation, SCM};
  \node[entry]  (a2) at (0,-0.55) {\S1.4 \ régularisation};
  \node[entry]  (a3) at (0,-1.10) {\S1.5 \ cas complexe};

  \node[ctitle] (t2) at (0,-2.05) {ch.\,2 --- Distributions};
  \node[entry]  (b0) at (0,-2.67) {\S2.1.1 gaussienne, Wishart};
  \node[entry]  (b1) at (0,-3.22) {\S2.1.2 elliptiques};
  \node[entry]  (b2) at (0,-3.77) {\S2.1.3 M-estimation};
  \node[entry]  (b3) at (0,-4.32) {\S2.2 \ cas complexe};

  \node[ctitle] (t3) at (0,-5.27) {ch.\,3 --- Riemann};
  \node[entry]  (c1) at (0,-5.89) {\S3.1 \ variétés, transport};
  \node[entry]  (c2) at (0,-6.44) {\S3.2 \ cône SPD (AI, LE)};
  \node[entry]  (c3) at (0,-6.99) {\S3.3 \ moyennes de Fréchet};
  \node[entry]  (c4) at (0,-7.54) {\S3.4 \ matrices orthog.};
  \node[entry]  (c5) at (0,-8.09) {\S3.5 \ optimisation};

  \node[ctitle] (t4) at (0,-9.04) {ch.\,4 --- Géométrie de l'information};
  \node[entry]  (d1) at (0,-9.66) {\S4.1 \ métrique de Fisher};
  \node[entry]  (d2) at (0,-10.21) {\S4.2 \ modèle gaussien};
  \node[entry]  (d3) at (0,-10.76) {\S4.5 \ BCR intrinsèque};
  \node[entry]  (d4) at (0,-11.31) {\S4.6 \ divergences};

  \coordinate (r) at (4.62,0); % garantit une largeur commune aux quatre cadres

  \begin{scope}[on background layer]
    \node[frame, fit=(t1)(a1)(a3)(r)]        (F1) {};
    \node[frame, fit=(t2)(b0)(b3)(r|-b0)]    (F2) {};
    \node[frame, fit=(t3)(c1)(c5)(r|-c1)]    (F3) {};
    \node[frame, fit=(t4)(d1)(d4)(r|-d1)]    (F4) {};
  \end{scope}

  %--------------------------------------------------------------- destinations
  \node[dest, draw=lcinq, text=lcinq, fill=lcinq!5] (ch5) at (10.9,-2.10)
    {\textbf{\textsc{iii}.1}\\Test d'hypothèses robustes\\\textit{sonar, RSO}};
  \node[dest, draw=lsix, text=lsix, fill=lsix!5] (ch6) at (10.9,-4.35)
    {\textbf{\textsc{iii}.2}\\Apprendre sur\\des covariances};
  \node[dest, draw=lsept, text=lsept, fill=lsept!5] (ch7) at (10.9,-6.90)
    {\textbf{\textsc{iii}.3}\\Apprentissage profond\\sur la variété SPD};
  \node[dest, draw=lhuit, dashed, text=lhuit] (ch8a) at (10.9,-9.55)
    {\textbf{\textsc{iii}.4}, \S2\\Transformers SPD};
  \node[dest, draw=lhuit, dashed, text=lhuit] (ch8b) at (10.9,-11.05)
    {\textbf{\textsc{iii}.4}, \S3\\Phase géométrique};

  \node[font=\scriptsize\bfseries, text=lhuit, anchor=south west] at (9.45,-8.75)
    {Perspectives};
  \node[font=\scriptsize\bfseries, text=lcinq!80, anchor=south] at (10.9,-1.05)
    {Contributions};

  %--------------------------------------------------------------- aretes
  % ---- vers le chapitre III.1
  \begin{scope}[edge, lcinq]
    \draw (F1.east) to[out=0,in=165]
      node[elab,pos=0.52,text=lcinq]{\S1.3 SCM \textperiodcentered\ \S1.4 régularisation\\\S1.5 circularité}
      (ch5.west);
    \draw ($(F2.east)+(0,0.30)$) to[out=0,in=195]
      node[elab,pos=0.62,text=lcinq]{\S2.1.2 SIRV\\\S2.1.3 Tyler, $M$-estimation}
      (ch5.west);
    \draw ($(F3.east)+(0,0.95)$) to[out=25,in=215]
      node[elab,pos=0.32,yshift=-13pt,text=lcinq]{\S3.2 distance affine invariante\\\S3.3 moyennes \textperiodcentered\ \S3.5 descente}
      (ch5.south);
    % contournement par la droite: ne pas traverser le bloc III.2
    \draw[rounded corners=4pt] ($(F4.east)+(0,0.55)$)
      -- (9.05,-8.49) -- (9.05,-2.90)
      node[elab,pos=0.86,text=lcinq]{\S4.1 test de Rao\\\S4.6 divergences}
      -- ($(ch5.west)+(0,-0.45)$);
  \end{scope}

  % ---- vers le chapitre III.2 (apprentissage sur des covariances)
  % Trois aretes seulement: ce chapitre n'emprunte rien au chapitre 4, et les
  % deux notions dont il a besoin en propre (modeles graphiques, matrices
  % aleatoires) sont introduites dans le chapitre lui-meme.
  \begin{scope}[edge, lsix]
    % ch.1: le regime dimensionnel, dont la theorie des matrices aleatoires est
    % la description fine.
    \draw ($(F1.east)+(0,-0.75)$) to[out=-25,in=170]
      node[elab,pos=0.25,text=lsix]{\S1.4 régime\\dimensionnel}
      ($(ch6.west)+(0,0.30)$);
    % ch.2: la vraisemblance elliptique penalisee de l'apprentissage de graphes.
    \draw ($(F2.east)+(0,-0.95)$) to[out=-12,in=178]
      node[elab,pos=0.20,yshift=-3pt,text=lsix]{\S2.1.2 elliptiques\\\S2.1.3 $M$-estimation}
      (ch6.west);
    % ch.3: cone, moyennes et descente riemannienne, communs aux deux travaux.
    \draw ($(F3.east)+(0,1.30)$) to[out=18,in=200]
      node[elab,pos=0.74,yshift=11pt,text=lsix]{\S3.2 cône \textperiodcentered\ \S3.3 moyennes\\\S3.5 optimisation}
      ($(ch6.west)+(0,-0.32)$);
  \end{scope}

  % ---- vers le chapitre III.3 (apprentissage profond sur la variete SPD)
  % Quatre aretes: c'est le seul chapitre du memoire a mobiliser les quatre
  % chapitres de l'etat des lieux.
  \begin{scope}[edge, lsept]
    % ch.1: la covariance apprise est une SCM, et son conditionnement decide
    % du seuil de ReEig. Contournement par le haut et par la droite.
    \draw[rounded corners=4pt] ($(F1.east)+(0,0.20)$)
      -- (12.85,0.20)
      node[elab,pos=0.55,text=lsept]{\S1.3 covariance empirique\\\S1.4 conditionnement}
      -- (12.85,-6.90) -- (ch7.east);
    % ch.2: Wishart et Wishart inverse donnent la lecture statistique des
    % moyennes arithmetique et harmonique.
    \draw ($(F2.east)+(0,-0.45)$) to[out=-20,in=150]
      node[elab,pos=0.82,text=lsept]{\S2.1.1 Wishart} ($(ch7.west)+(0,0.42)$);
    % ch.3: le gros de la dependance.
    \draw ($(F3.east)+(0,0.05)$) to[out=0,in=180]
      node[elab,pos=0.48,text=lsept]{\S3.2 cône (AI, LE) \textperiodcentered\ \S3.3 moyennes\\\S3.4 Stiefel \textperiodcentered\ \S3.5 rétractions}
      (ch7.west);
    % ch.4: c'est lui qui fait de la liste des moyennes une famille indexee par
    % un modele statistique.
    \draw ($(F4.east)+(0,-0.35)$) to[out=10,in=250]
      node[elab,pos=0.30,text=lsept]{\S4.2 distance de Rao gaussienne\\\S4.6 divergences}
      ($(ch7.south)+(0.35,0)$);
  \end{scope}

  % ---- vers les perspectives
  \begin{scope}[edge, lhuit, dashed]
    \draw ($(F3.east)+(0,-0.85)$) to[out=-15,in=180]
      node[elab,pos=0.62,text=lhuit]{\S3.3 moyenne de\\Karcher pondérée} (ch8a.west);
    \draw ($(F3.east)+(0,-1.35)$) to[out=-30,in=180]
      node[elab,pos=0.82,text=lhuit]{\S3.1 transport\\parallèle} (ch8b.west);
    \draw[rounded corners=4pt] ($(F1.east)+(0,-0.40)$)
      -- (5.25,-0.95) -- (5.25,-12.05)
      -- (10.9,-12.05)
      node[elab,pos=0.88,text=lhuit]{\S1.5 pseudo-covariance} (ch8b.south);
    \draw ($(ch7.south east)+(-0.90,0)$) -- ($(ch8a.north east)+(-0.90,0)$);
  \end{scope}

  %--------------------------------------------------------------- legende
  \node[anchor=north west, font=\tiny, text=lctx, align=left] at (0,-12.40)
    {\textcolor{lcinq}{\rule[0.4ex]{7pt}{0.5pt}}~outil mobilisé en \textsc{iii}.1
     \quad
     \textcolor{lsix}{\rule[0.4ex]{7pt}{0.5pt}}~en \textsc{iii}.2
     \quad
     \textcolor{lsept}{\rule[0.4ex]{7pt}{0.5pt}}~en \textsc{iii}.3
     \quad
     \textcolor{lhuit}{\rule[0.4ex]{3pt}{0.5pt}\hspace{1.5pt}\rule[0.4ex]{3pt}{0.5pt}}~piste ouverte en \textsc{iii}.4
     \\
     (les numéros de chapitre repartent à 1 dans chaque partie)};
\end{tikzpicture}
